% -*- coding: utf-8 -*-
% !TEX program = xelatex
\documentclass[plain,noanswer,medmath]{../../template/jnuexam}
\usepackage{../../template/kysx}

\begin{document}


\examtitle{1987}{5}


\loadproblems{4}{1987P4}%载入试卷四


\makepart{判断题}{本题满分10分，每小题2分}


\useproblem{4}{1}{1}%【同试卷四第一(1)题】


\useproblem{4}{1}{2}%【同试卷四第一(2)题】


\begin{problem}
若函数$f(x)$在区间$(a,b)$内严格单调增加，
则对于区间$(a,b)$内的任何一点$x$有$f'(x)>0$．\tickout{}
\end{problem}


\begin{problem}
若$A$为$n$阶方阵，$k$为常数，%而$|A|$和$|kA|$为矩阵$A$和$kA$的行列式，
则$|kA| = k|A|$．\tickout{}
\end{problem}


\useproblem{4}{1}{5}%【同试卷四第一(5)题】


\makepart{选择题}{本题满分10分，每小题2分}


\begin{problem}
函数 \pickin{} 在其定义域内连续．
\begin{abcd}
\item $f(x) = \frac{1}{x}$
\item $f(x) = \begin{cases}
  \sin x,& x \le 0; \\
  \cos x,& x > 0
\end{cases}$
\item $f(x) = \begin{cases}
  x + 1, & x < 0; \\
      0, & x = 0; \\
  x - 1, & x > 0
\end{cases}$
\item $f(x) = \begin{cases}
  \frac{1}{|x|}, & x \ne 0; \\
              0, & x = 0
\end{cases}$
\end{abcd}
\end{problem}


\useproblem{4}{2}{2}%【同试卷四第二(2)题】


\useproblem{4}{2}{3}%【同试卷四第二(3)题】


\useproblem{4}{2}{4}%【同试卷四第二(4)题】


\begin{problem}
对于任意二事件$A$和$B$，有$P(A-B)=$\pickout{}
\begin{abcd}
\item $P(A) - P(B)$．
\item $P(A) - P(B) + P(AB)$．
\item $P(A) - P(AB)$．
\item $P(A) + P(\bar B) - P(A\bar B)$．
\end{abcd}
\end{problem}


\makepart{计算题}{本题满分20分，每小题4分}


\begin{problem}
求极限$\lim_{x \to +\infty} \frac{\ln \left(1+\frac{1}{x}\right)}{\arctan x}$．
\end{problem}


\useproblem{4}{3}{2}%【同试卷四第三(2)题】


\useproblem{4}{3}{3}%【同试卷四第三(3)题】


\begin{problem}
计算定积分$\int_{\frac{1}{2}}^1\e^{\sqrt{2x-1}} \dx$．
\end{problem}


\begin{problem}
求不定积分$\int \frac{x\dx}{x^4 + 2x^2 + 5}$．
\end{problem}


\makepart{}{本题满分10分}


\begin{problem}

考虑函数$y=x^2$, $0 \le x \le 1$（如图），问：\par
\centerline{\includegraphics[width=0.5\linewidth]{1987E}}
\begin{enumerate}
\item $t$取何值时，图中阴影部分的面积$S_1$与$S_2$之和$S = S_1 + S_2$最小？
\item $t$取何值时，面积之和$S = S_1 + S_2$最大？
\end{enumerate}
\end{problem}


\makepart{}{本题满分5分}%【同试卷四第六题】

\useproblem{4}{6}{0}%


\makepart{}{本题满分8分}

\begin{problem}
假设某产品的总成本函数为$C(x) = 400 + 3x + \frac{1}{2}x^2$，
而需求函数为$p=\frac{100}{\sqrt{x}}$，其中$x$为产量（假定等于需求量），$p$为价格．试求：\par
\begin{enumerate*}
\item 边际成本；
\item 边际效益；
\item 边际利润：
\item 收益的价格弹性．
\end{enumerate*}
\end{problem}


\makepart{}{本题满分8分}%【同试卷四第八题】

\useproblem{4}{8}{0}%


\makepart{}{本题满分7分}%【同试卷四第九题】

\useproblem{4}{9}{0}%


\makepart{}{本题满分6分}%【同试卷四第十题】

\useproblem{4}{10}{0}%


\makepart{}{本题满分8分}

\begin{problem}
已知离散型随机变量$X$的概率分布为：
$$P\{X=1\} = 0.2,\; P\{X=2\} = 0.3,\; P\{X=3\} = 0.5.$$
\begin{enumerate*}
\item 写出$X$的分布函数$F(x)$；
\item 求$X$的数学期望和方差．
\end{enumerate*}
\end{problem}


\makepart{}{本题满分8分}%【同试卷四第十二题】

\useproblem{4}{12}{1}%


\end{document}
